\section{Aztecs and Human Sacrifice}

On the supposition that Greenland was larger and connected the Eastern and Western hemispheres, these Hyperboreans, related to the Cro-Magnon man, would have wandered as far as Meso-America. As evidence of a connection, Evola points to the Aztec claim to be from Atzlan (the land of light) and the Toltecs to Thule. Presumably, these migrations were mixed with Orientals and the Lemurians. Evola seems to waver on the solarity of the Aztecs. In Pagan Imperialism, he explains:

\begin{quotex}
from China to Greece, from Rome to the primordial Nordic groups, then up to the Aztecs and the Incas, nobility was not characterised by the simple fact of having ancestors, but by the fact that the ancestors of the nobility were divine, unlike those of the plebeians, and to which though it can remain faithful, also through the integrity of blood (in the caste system, the principle of heredity was valid not only for the higher castes, but also for the lower ones). The nobles originated from ``demigods'', that is to say, from beings who had actually followed a transcendent form of life, forming the origin of a tradition in the higher sense, transmitting to their lineage a blood made divine, and, along with it, rites, that is, determinate operations, whose secret every noble family preserved, which allowed their descendants to continue the spiritual conquest from where it had previously reached, and to lead it gradually from the virtual to the actual.
\end{quotex}


Thus ``family values'' do not refer to the bourgeois nuclear family as commonly assumed today. The Vedas show that the ``family'' encompasses generations.

\begin{quotex}
Thus, from the traditional point of view, not-having-ancestors distinguishes the plebeian from the patrician less than not-having-rites. In Aryan hierarchies, a single characteristic differentiated the higher castes from the lower: rebirth. The arya, as opposed to the shudra (the one who serves), was the dvija, the twice-born. The assertion of the Manavadharmashastra (II, 172), that the brahman himself, if he left out initiation, would no longer be differentiated from the one who serves, the shudra, is indicative.
\end{quotex}


Keep in mind here, that for Evola, the shudra are of a different race from the Brahmans, at least in terms of ``race of the spirit''. As for the Nordics, Evola writes:

\begin{quotex}
The Nordic nobles were noble because, in their blood, they carried the blood of the Asen, of the ``celestial'' forces in continuous struggle with the elemental beings. The nobility of the great medieval orders of chivalry --- among which the most significant were the Templars --- was also tied to initiation. One of the weakest points in Nietzsche's conception is precisely his biological naturalism, which, in most cases, diminishes and secularises his aristocratic idea, carrying it down to the level of the ``blond beast''.
\end{quotex}


Here we see again Evola's emphasis on the greater battle, waged within. As for the Aztecs, Evola expressed a somewhat different opinion in Revolt against the Modern World, where he seems to regard them as a decadent form, but anyone can read that, so I won't repeat. For a rather different take on this question, I'll refer to \textit{The Battle for Amerindia} by the Traditionalist Catholic \textbf{Solange Hertz}.

\begin{quotex}
so-called ``underdeveloped'' peoples languishing in the lower stages of evolution have actually degenerated from higher cultures.
\end{quotex}


In one sense, she is referring to the Amerindians of Mesoamerica, although, in her view, they were not nearly as degenerate as modern Europeans who have lost connection with their own past. This is a sobering thought which should not be passed over lightly. To regard whites as somehow superior (in culture, creativity, productivity, IQ, and so on) and thus immune to the fate of other peoples, misses the point. Most of the literature of that genre focuses on inferior peoples being admitted into formerly white lands. Mrs. Hertz is claiming that perhaps it is the whites who, having abandoned their spiritual heritage in the pursuit of decadence and economic gain, are in truth the spiritually degenerate people. As for the sacrificial practices of the Aztecs, she offers this historical tidbit:

\begin{quotex}
Las Casas the Bishop of Chiapas even went so far as to maintain that the Spaniards had no right to interfere forcibly with the human sacrifices offered by pagan Amerindians, detecting in these, as would Joseph de Maistre two hundred years later, a perverted but profound religious instinct which did not flinch from surrendering to their deities what they deemed most precious!
\end{quotex}


\textbf{Joseph de Maistre}, by the way, was an author esteemed by Evola.

As a last comment which can either provide food for thought or be ignored totally, we can note, as did Evola, that Mesoamerica was dedicated to the Serpent. Mrs Hertz tells us that the Lady of Guadalupe who appeared to Juan Diego is a transliteration of the Aztec word \emph{Coatlalopeuh}, which means the ``Crusher of the Serpent''.

At the very least, Mrs. Hertz provides a more nuanced view of the relationship between the Spanish and the Aztecs that what we hear from the hard left every \textbf{Columbus Day}. Bishop Las Casas was clearly against forced conversions and was able to see what was best in the rites of the pagan Aztecs.


\flrightit{Posted on 2011-10-16 by Aeneas}

\begin{center}* * *\end{center}

\begin{footnotesize}
\begin{sffamily}
\texttt{Will on 2011-10-17 at 07:53 said:}

Excellent post, but I don't remember Evola's view of the Aztecs being any different in Revolt, so I went back and checked. On pages 232-3, he says that the Aztecs, Toltecs, and Incas were the bearers of solar civilization, overcoming the Mayas and other telluric peoples. Historically this would seem to parallel the Aryan invasion of India and the Dorian invasion of Greece. However, by the time of the European conquest, these civilizations had indeed degenerated, and that is the decadent form to which your post refers. Analogously, the Spartans of the 3rd century B.C. were a far cry from those of the 6th.

The Catholics were not the first to `crush the serpent' in Mesoamerica. The Aztecs built their capital Tenochitlan (now Mexico City) at the sight of an eagle holding a snake in its claws. The eagle is a traditional Olympian symbol, and here it would seem to symbolize the conquest of the telluric by the solar.

\hfill

\texttt{EXIT on 2011-10-17 at 10:13 said:}

A.:''Maybe WN groups are seeking out initiations privately.''

I'm no WN, but I don't understand how you can talk about initiation when gornahoor supports self-initiation (which Guenon and co. referred to as pseudo-initiation); hence gornahoor is in no better position than WN.

\hfill

\texttt{James O'Meara on 2011-10-17 at 19:20 said:}

Still pasting Evola quotes into content-free arguments, I see.

WNs who are... ? 

``Their goals... Which are... ? Have you ever, at any time on this blog, ever, ever named a ``WN'' or ``their goals''? Of course not, then someone would be able to evaluate your arguments.

You could substitute ``Americans'' or ``pawnbrokers'' or ``these damn kids with their rock and roll'' or whatever and make as much a contribution to thought.

You're like one of those stuck-up Euros who deplored America's ``racism'' from safely White communities, until the Arabs started burning police cars in Paris and committing EVERY rape in Stockholm true fact! If you think America will have any `initiatory rites' or `esoteric Christianity' when the whole thing looks like Detroit, you're crazy.

If you think Evola would support importing millions of North Africans into Italy, you're crazier. 

[

Although we never implied that, it follows the logic of the Left Hand Path. Choose the more difficult experience.

As for the program of the WNs, we've found someone to discuss that and it will appear on the Right Hand Path blog directly.

C!

]

\hfill

\texttt{James O'Meara on 2011-10-17 at 19:46 said:}

Also, it's pretty rich for you to quote from Pagan Imperialism, a book that Evola wrote to suck up to the Fascisti, and then re-wrote, dropping the Roman stuff and lauding the Nordics, to suck up to the Nazis even, later, swanning around in a Gestapo uniform , in order to critique ``White Nationalists.'' 

[Yes, James, life is full of such contradictions. The trick is to foresee who to suck up to.

Glad to have you aboard again. C.]

\hfill

\texttt{Cologero on 2011-10-17 at 20:31 said:}

To Exit:

Besides the Traditional Catholic initiations, whose validity is not accepted unanimously by the early masters of Tradition, I have been initiated by a Tibetan lama. I will make the documentation available online in the near future.

Cologero


\hfill

\texttt{Cologero on 2011-10-18 at 00:45 said:}

Re Initiations

Besides initiations into Traditional Catholic and Buddhist forms, I have other initiations, whose regularity may be doubted by some.

The truth of the matter is that in the Kali Yuga much more can be accomplished by a man who knows how. As Jesus told us, ``Be as innocent as the Holy Spirit, but as cunning as Lucifer.'' Heaven must be taken by storm in our age, a man must be clever and find his way by piecing things together.

The greater struggle is more difficult in our time. However, winning that battle will give us more strength, power and will than it would in more pacific times. Negativity and chaos can and must be transmuted into order and transcendence. The greater the negativity, the greater the possibilities of overcoming.

So, Exit, if you are expecting someone to take you by the hand like some kindergarten student and lead you into some perfect initiation, then you are not the type of man that Gornahoor is intended to appeal to.

For anyone interested, these are those groups; however, \emph{a chacun son gout}. Whatever their demerits, they still teach experiential practice. By studying the more traditional practices, one learns to see what they are really saying. Furthermore, by tying personal experience in with a Traditional metaphysical teaching, they each reinforce each other. Given our situation today, it is what we have.

Transcendental Meditation

Gurdjieff Fourth Way work

Kaballah

Rune Guild (on my own as I never took formal initiation)

\hfill

\texttt{Matt on 2011-10-18 at 02:31 said:}

Looking forward to the future documentation on the initiation by the Lama. Did it take place in the Himalayan region?

\hfill

\texttt{Charlotte Cowell on 2011-10-18 at 17:52 said:}

This is tangential to your post on Aztecs, but the comment about initiation and rebirth caught my eye. I am presuming that most people will agree that true initiation is not conferred in a mundane ritualistic way by one human being upon another, as may happen in a `secret society'. The right hand, after all, is not meant to know what the left is doing in order for such divine magic to occur. True initiation is (in my view at least), something that takes place via the vertical hierarchy of angelic forces and, ultimately, the Holy Spirit. (The identity of the initial `psychopomp' is perhaps an absolute mystery, something which defies explanation and description, characterised simply by `its' function). Most people would also agree, however, that as no man is an island, we are all dependent upon one another, to some extent, for the catalytic energy that might assist with our personal development. When it comes to the matter of rebirth in relation to initiation one must necessarily be dependent upon an awakening of memory -- for how can one be conscious of past lives without remembrance. So I put the question to you, are we therefore dependent upon others (our `karmic relations' if I might get Steinerish for a moment) for this reawakening? Even more -- given the ultimately androgynous condition of the complete human entity, reflection of the divine -- is it essential to encounter our `soul mate' or `twin flame' (or however one must define such beings and aspects of our self) in order to take certain initiatory steps forward? A Hindu or Buddhist might say no, it is possible to achieve this state by individual effort alone (albeit with intense scriptural meditation, which might possibly be seen as a form of programming), whereas a Kabablist might say yes... I wonder what the `acceptable' Christian position is on this given that reincarnation is still very much taboo (probably with very good reason) for the exoteric church... .

\hfill

\texttt{Perennial on 2011-10-19 at 03:47 said:}

Charlotte, reincarnation, as it is generally understood, is taboo amongst the Traditionalists as well, denounced by both Evola and Guenon. They did believe in a form of ``reincarnation'', however. So when you say reincarnation, what do you mean by that?

\hfill

\texttt{Charlotte Cowell on 2011-10-19 at 12:18 said:}

well Perennial, doubtless it's an immense mystery. One thing is for sure, exploration of the reincarnation idea can lead one right to the edge and beyond of the abyss, which is surely why it's taboo. A necessary part of pathworking? Perhaps, but fraught with spiritual, psychological and emotional danger if one falls into the Qlipoth, as human beings almost invariably will (must, given our imperfect state?). But I would still say that one should not `go there' unless utterly compelled in some unavoidable way. ie, if one's soul is trapped in the `eighth sphere' without hope of return unless one goes back to the beginning and starts again... maybe that's all of us at some stage or other.

But looking at it another way, the body is a temple for the soul/spirit, so who knows what and/or whom might walk through those doors, invited or uninvited? If you or I achieve a certain state of receptivity whereby our `selves' (our ordinary ego selves) are turned outwards, a space is made within. In this context, one might experience many, many lives during one life on Earth. I might say without trace of irony, I've lived but once -- this is the first and last -- yet I'm not only Eve but Artemis, not only the Magdalene but the Jezebel. Marie Antoinnette and sea priestess, every woman in fact, while you are every man from Adam to Apollo, Jesus Christ and Caesar, and that's without even considering the animal kingdom. Is that reincarnation enough, or is there more? While I feel this is the acceptable and Christian line and also in fact the simple truth, there is more when we go through the looking glass (which perhaps we should not do. Probably we should not).

For many years following my conversion to Christianity I simply didn't think about reincarnation, I really didn't care if it was a fact or not, I considered it to be one of those things that one need not think about (along with Atlantis and Shambhala and all kinds of other `grey areas' I was led to face in recent years). What made the difference for me is that after many years of what I'd describe as a dark night of the soul -- a true lonely wilderness -- I met another of those people that come into one's life from time to time and trigger a massive change. I have been in love before, but when I met this particularly person I was forcibly and instantaneously struck by the fact that we'd known each other forever, which is not something I'd ever experienced before.

We saw each other only very rarely and briefly over the next 3 years or so, but at some point this (very materialistic down to earth) man said he felt like our relationship was magical, as if we'd known each other in a past life. This sounds very romantic and even slushy no doubt, but perhaps you can appreciate that for someone with a very traditionally Christian mindset this was an idea with revelatory power that opened up doorways (real or imaginary) that I never knew existed (he'd said what I had already been thinking for a long time, though I wrestled with the idea long and hard, wondering if it was just an illusion. As it happens, I probably did fall into the domain of Maya, but we all go through this so I'm not afraid to say it).

But even this is a simplification given that the Imperial task is to be fully human, with all the human loves and losses that entails. I am struggling a bit to define how my mindset began to change -- or perhaps, how my perspective was altered -- but my continual longing to be with this person had the effect of leading me to look backwards as far as possible to the start of time, as I explored the idea of soulmates and the originally unified (androgynous) being (human and divine).

Some aspects of this exploration posed me with no difficulties while others appeared as insurmountable problems. The quality of time, for instance, is not something that phases me, as from the very start of my esoteric spiritual journey I've been fully aware (and in no doubt) that it has no meaning beyond this dimension. It binds us `here', but we are entirely free of it `there', as it is for space, the other side of the same continuum. So in this respect `looking back' has to be analogous with `looking in', it's simply a matter of evoking the state one wishes to explore -- ancient Greece for instance, or Genesis. (Of course one could not accomplish this without recourse to myth for our archetypal references).

So I had little difficulty drawing a conclusion about my `past life' in ancient times with this individual, because whether or not we as personalities really were reborn again and again until the present age, the analogy still stood at an archetypal level. (even if it is only true in one's own mind, it is still a tool one might usefully use when navigating the road of return. We are born male and female together, and needs must comprehend our opposite side in order to get back to where we were. Whether another human can only ever be a `best possible match' of anima or animus, or whether (as kabbalists would say), they truly are another half to the self, is probably a matter of opinion so far as the world today is concerned.

In the context of physics -- quantum mechanics, superstring theory and so on -- the Law of Vibration described by the Kybalion makes the association simple. Whatever really happens to the soul throughout earth ages, the resonance is real, the analogy works, the echo resounds and time/space are traversed, even escaped in an eternal sense. Reincarnation is almost a red herring in this context, as everything is already and always there in the realm of eternity, we only need remember it and find it within ourselves. Without going mad, losing our faith or killing ourselves/ each other :-)

Perhaps I am not explaining myself very well as there was so much emotional content to the experience (overwhelmingly so), but suffice it to say that this one person in particular fully awakened my memory, leading me to believe that we had occupied a similar position in relation to one another at another point in the space-time continuum. Possibly more than once. So there is a powerful soul-mate element to all of this and again I think it's very Steinerish in flavour (not deliberately), in that the idea of reincarnation seems to be inextricably linked to karmic relationships.

So to answer your question, in the sense pertaining to Maya it is `all in the mind', but in a concrete spiritual sense we are eternal and therefore contain the whole of time and space. If a certain note is struck within our souls it will vibrate all the way along that line from start to finish ad infinitum. it's a matter of comprehending the essence -- the note -- that corresponds to a certain archetype, past or future.

As an addendum, it's worth noting that some individuals under hypnosis have come out with detailed information from the past that only appears to be explicable in the context of either reincarnation or mediumistic practice.

\hfill

\texttt{EXIT on 2011-10-21 at 18:26 said:}

C. wrote, ``I have been initiated by a Tibetan lama.''

Is this the same fictional Tibetan adept from whom Blavatsky attempted to derive her authority? Speaking of which the ``society of independents'' sounds quite similar to Encausse's ``independent group for esoteric studies''. But even if I give you the benefit of the doubt that you did get a Buddhist initiation, it is not as if one receives a simple initiation and automatically attains enlightenment. It takes years and years with a guru before one is even in a position to attain such, that I would not take anyone's word that he had firsthand knowledge of metaphysical states at face value. Seeing as though you do not practice Buddhism outwardly, at least not in your writings, I think it safe to say that this initiation was not that serious.

\hfill

\texttt{Cologero on 2011-10-22 at 16:51 said:}

Exit, you are a special kind of idiot so dear to me, the \emph{idiot savant}, someone who knows how to do all the calculations yet does not quite understand what they are for. You can start with what we wrote on \textit{Madhyamika}\footnote{\url{http://www.gornahoor.net/?p=299}}. An effete display of erudition is of no benefit to anyone; our intent here is to outline the lineaments of a possible future tradition, which will not be created by me, yet it will happen. The first step, as in the Madhyamika, is to overcome all points of view; without that, there is no point in proceeding. Hence, we have related it to many posts, particularly on the Evolean idea of the First Trial. Interested readers may also look up our posts on Josephin Peladan and see how they relate. It is important to read in depth, not just horizontally.

\hfill

\texttt{Charlotte Cowell on 2011-10-22 at 17:39 said:}

To C! and Exit, firstly, I'd avidly love to hear more about the Tibetan initiation.

Secondly, my own view on initiation is that it is conferred by the Holy Spirit and one is able to perceive the presence of a guiding entity. Some might define this as the guardian angel, or as a psychopomp. Another that should probably be said is that it (initiation) comes in stages and one had most definitely better be ready, willing and able at all times to drop absolutely everything when required to take the next step.

Next, more pertinent to this particular point, anyone on the initiatory path -- or the `royal road' as a Sufi might say -- is (from what I can tell, based on the experience of every true initiate I've ever known of), going to head towards India/Tibet in more ways than one. It's a simple matter of it being what happens.

You are right to point out the years of discipline required by way of preparation -- it's easy to see how Buddha laid the carpet down for Christ in the direction of bodily control, right motivation and so on. But this is not all there is to it. The wisdom of the ages is to be found in all corners of the Earth and the Himalayan teaching is a one of these -- gold is produced in this region, remember?

Also remember that the teachers of this region are masters in the art of telepathic communication, astral and etheric travel. Spend enough time in those realms and one will encounter Tibetan guides. Seek and ye SHALL find, but how ardently one must seek... ..well, it has to be longed for above all other things and override every other commitment in life.

\hfill

\texttt{Charlotte Cowell on 2011-10-22 at 17:52 said:}

Exit, this comment of yours:

``it is not as if one receives a simple initiation and automatically attains enlightenment''

Well as it happens, these things DO occur in the twinkling of an eye and surely one will be amazed in that moment of eternity! Indeed, it will be a moment of such import that the impression it makes on the soul will be powerful enough to sustain that being through the trials that follow. It WILL constitute the defining moment of one's life. The light itself must be strong enough to illumine the way for many long and otherwise dark years, without it one would simply give up in despair.

Very specific to what is being described as a Buddhist initiation, from what I know I can say that it's certainly possibly to receive an instantaneous `dharma transmission', the efficiacy of which depends on both the capacity to `pass over' such knowledge on the part of the `carrier', and the twin capacity to receive on the part of the neophyte, or pupil. Trust and purity of intention are fundamental to the success of this procedure and, I would suggest, an element of conducive personal chemistry is also required for such face-to-face, heart-to-heart, spirit-soul contact.

So yes, believe it or not, one IS transformed in a flash of light, praise be to God!

\hfill

\texttt{EXIT on 2011-10-23 at 09:32 said:}

CC:''So yes, believe it or not, one IS transformed in a flash of light''

If you believe that you're only fooling yourself. How about Zen satori? 

\hfill

\texttt{Charlotte Cowell on 2011-10-24 at 05:38 said:}

`There are two ways to be fooled. One is to believe what isn't true; the other is to refuse to believe what is true' (Kierkegaard)

It was Jung who said that he didn't have to believe any more because he KNOWS... .


\hfill

\end{sffamily}
\end{footnotesize}

